%!TEX root = main.tex

In this section, we describe the main ideas used in the formalization of HiFIRRTL.

Let us focus on a HiFIRRTL module, say $M$.

A structure called circuit diagram is used in the formalization to represent HiFIRRTL programs.

\begin{definition}
A \emph{circuit diagram} $D$ is a triple $(ss, ce, cm)$, where $ss$ is a sequence of statements, $ce$ is a component environment that maps each component to its type, and $cm$ is a connect mapping that maps each component to an expression that is connected to it. 
\end{definition}

A HiFIRRTL module $M$ will be first transformed into a circuit diagram $D_M = (ss, ce, cm)$, where $ss$ is the sequence of statements in $M$, $ce$ is the component environment obtained from the variable declarations in $M$, possibly with the types of some variables missing or partial, and $cm$ is the trivial mapping that maps each component to $\bot$ (Intuitively, this means that the connects of the components have not be determined yet). The missing or partial  types in $ce$ will be completed by the compilation passes during the transformation of $M$ into a LoFIRRTL module. Moreover, the connect mapping $cm$ will be updated by the ExpandWhens compilation pass. Furthermore, the statement sequence $ss$ will be updated by ExpandConnects, ExpandWhens, LowerTypes, and RemoveReset, while they keep unchanged during the other compilation passes.  

The formal semantics of HiFIRRTL will be defined by the axiomatic semantics of the compilation passes, plus the operational semantics of LoFIRRTL.
%InferTypes, InferWidths, InferResets, ExpandConnects, ExpandWhens, LowerTypes, and RemoveReset.  
%
More precisely, the axiomatic semantics of each compilation pass $\pi$ will be specified by a formula $\varphi_{\pi}(\cdots)$, and 
%Note that the sequence of statements in $ss$ will keep unchanged in the compilation passes.
%
%ce1 = φ_InferType(ss) /\
%ce2 = φ_InferWidth(ss, ce1) /\
%ce3 = φ_InferReset(ss, ce2) /\
%(ss1, ce4) = φ_ExpandConnect(ss, ce3) /\
%(ss2, cm) = φ_ExpandWhen(ss1, ce4) /\
%(ss3, ce5, cm1) = φ_LowerTypes(ss2, ce4, cm) /\
%(ss4, ce6, cm2) = φ_RemoveResets(ss3, ce5, cm1)
the semantics of $M$ is defined by the formula
\[
\begin{array}{l c l}
\varphi_M & \eqdef & \varphi_{ResolveKinds}(ss, ce_1)  \wedge  \varphi_{InferTypes}(ss, ce_1, ce_2)\ \wedge \\
&& \varphi_{InferWidths}(ss, ce_2, ce_3) \wedge  \varphi_{InferResets}(ss, ce_3, ce_4)\ \wedge \\
&& \varphi_{ExpandConnects}(ss, ce_4, ss_1, ce_5) \wedge \varphi_{ExpandWhens}(ss_1, ce_5, ss_2, cm)\ \wedge \\
&& \varphi_{LowerTypes}(ss_2, ce_5, cm, ss_3, ce_6, cm_1) \wedge  \varphi_{RemoveReset}(ss_3, ce_6, cm_1, ss_4, ce_7, cm_2).
\end{array}
\]
Note that each formula $\varphi_{\pi}$ actually specifies the graph of a function. For instance, the formula $\varphi_{InferWidths}(ss, ce_2, ce_3)$ specifies a function that takes $ss$ and $ce_2$ as inputs and has $ce_3$ as the ouput.  The final output $ss_4$ is actually a sequence of statements already in LoFIRRTL form, whose operational semantics has been defined in Section~\ref{sec-lofirrtl-sem}.

In the sequel, we shall illustrate the definition of the semantics of compilation passes by describing $\varphi_{InferWidth}$ and $\varphi_{ExpandWhen}$. The semantics of the other compilation passes can be found in Appendix~\ref{sec-appen-sem-cpass}.  

%Let us fix $ss = (s_1, \cdots, s_n)$ in the sequel.


$\varphi_{InferWidths}$

$\varphi_{ExpandWhens}$
